\documentclass[a4paper,11pt]{article}

\usepackage{revy}
\usepackage[utf8]{inputenc}
\usepackage[T1]{fontenc}
\usepackage[danish]{babel}
\usepackage{amsmath,amssymb}

\revyname{MatRevy}
\revyyear{2026}
\version{2026-10-09}
\eta{4 minutter}

\title{Blokuge 8}
\author{Augusta '20}
\melody{DtMF - Bad Bunny}

\begin{document}
\maketitle

\begin{roles}
\role{S}[Sanger] Studerende
\end{roles}

\begin{song}
\sings{S}[Vers] 
Ah ah ah ah 
Jeg har aldrig vær't særlig god tidligt om morgen 
Direktørtid er mere min stil, se mig vågne lidt over 10 
Op på cyklen og hen at læs' til 
Min mundtlig' eksam' 
Min mundtlig' eksam'

Ruller ind på HCØ 
Hvorfor står folk i kø 
Hvad er dog denne ballad'
De fortæller mig alle lokaler er taget 
(\textit{"Du kommer for sent bro. Haha."})
Si'r du til mig at 12 det er sent eller hvad? 
(\textit{"Kom venner, vi går bare hjem og læser i stedet".})

\sings{S}[Vers]
Nej! 
Det et negativt mindset du har dig der  
Én ledig tavle må der da vær'
Kom med - ned af den svalegang
Mens jeg tjekker lokaler og synger min sang
A101 helt ned til Aud 10, åh nej, jeg går lig' forbi 
For hvert en'ste lokal' er der mennesker i

Og jeg sender et dræberblik ind gennem vinduet 
Til folk som bar' tager et lokal' for at sid' ned 
Og blokerer os andre fra tavlen og kridtet 

Er for B-menns'k til oplæsningsugen
(Men) der findes en ny dag imorgen 
Bar' vent!

\sings{S}[Omkvæd] 
For jeg går hjem på værelset og stiller mit væk'ur 
Til sygt tidligt om morg'n, nærmest nat eller noget i den dur 
Ey, hvis jeg' op klokken 8, så skal det nok gå
Det jeg' sikkert den eneste der har tænkt på 

\sings{Alle}[Omkvæd] 
Så jeg går hjem på værelset og stiller mit væk'ur 
Til sygt tidligt om morg'n, nærmest nat eller noget i den dur 
Hvis jeg' op klokken 8, så skal det nok gå
Det jeg' sikkert den eneste der har tænkt på 

\sings{S}[Vers] 
Ey, går tilbag' ind og prøver aud 4 
Men der sidder sgu nogen
Tavlerne nede ved fysikerne
Ja der står også nogen
Selv begge tavler i snullet er taget, der står også nogen
Der er altid nogen:(

Ey! En ting som mit stressniveau ikke kan tåle
Det er fag som har øvelsestimer der ligger pladask lig' midt i blokuge 8
Og de fylder det hel', det hel', det hel'
Jeg forstår slet ik' at man må det

I blokuge 8 (\textbf{A}: \textit{Blokuge 8})
I blokuge 8 (\textbf{A}: \textit{Blokuge 8})
Hørt' ik' min alarm, på uni igen i min morgenkåbe 
I blokuge 8 (\textbf{A}: \textit{Blokuge 8})
I blokuge 8 (\textbf{A}: \textit{Blokuge 8})
Skal fremlægge 12 beviser på tavlen, kan slet ikke nå det
Så jeg går hjem og


\sings{S}[Tale:] 
\textit{Uanset om jeg sætter min alarm til klokken 8, 7, 6, 5 eller 4 om morgenen, så lykkes det stadig de andre at være der før mig. 
Men jeg har en plan. Hvis jeg går ind om natten og sover i et A-lokale, så vil jeg være den første som er derinde om morgenen. 
Jeg hader blokuge 8. 
}

\sings{S}[Vers] 
Jeg går hjem på mit vær'lse, pakker en taske 
Med sov'pose, trangia, telt og vandflaske 
Vil campe på skolen og være den første  
Ey, hvis jeg' der aft'nen før, ja så skal det nok gå
Det må jeg vær' den eneste der har tænkt på 


\sings{Alle}[Omkvæd:] 
Vi kommer aftn'en før og vor's tasker har telt i  
så nu er der dømt camping i alle lokaler på uni 
At ku' øve på tavle, det kan du glemme 
Du må bare læse derhjemme 
(Jubel)
Tænk at du troede du vil' vær' her før os (ay, wuh)
\end{song}
\end{document}
